\section{Developing FishBot v0.7}
\label{sec:development}

This section reports the development phases: what the instrument was shown to be capable of
detecting, what an open adversary search found, and which candidate mechanisms were tested and
rejected. All measurements here are on \emph{training} material and are exploratory in the sense of
\S\ref{sec:evaluation-multiplicity}. None supports the paper's primary claim, and every pooled
interval reported in this section has a lower bound below the detection floor of
\vsevenPrimaryFloor{}~pp, so each is a measured effect rather than a confirmed one.

\subsection{The v0.6 frontier, and the choice of comparison target}
\label{sec:development-frontier}

\vsix{} deployed its policy with test-time search disabled, while also having measured that search
working. The configurations available at the end of that cycle therefore form a \emph{frontier}
rather than a point, running from the deployed blueprint to progressively more expensive searching
configurations.

This distinction drove the most consequential methodological choice in the programme. Beating the
deployed policy is the easy comparison, and several unfitted single-switch deviations of it already
do. \textbf{The comparison target was therefore fixed as \fcheap{}}, the cheapest configuration that
is genuinely on the frontier, and \S\ref{sec:design-protocol} records that the choice was registered
in advance. Table~\ref{tab:configs} defines every configuration this paper compares.

\subsection{What the instrument can detect}
\label{sec:development-instrument}

The detection-floor procedure of \S\ref{sec:evaluation-floor} was calibrated in this phase and
produced the two findings reported there: the floor is \vsevenPrimaryFloor{}~pp, and it does not
decrease with evaluation games. Two further results shape later sections.

\paragraph{Transcript inversion contracts the posterior, and is not an attack.} Because the target
team is deterministic and published, an adversary can invert its transcript, asking for each observed
ask which deals would have produced it and contracting the posterior accordingly. Measured, this
contracts the deal posterior by approximately \vsevenBitsPerAsk{} bits per observed ask beyond the
certificates the rules already supply. Converting those bits into play is worth
\vsevenInversionUnfitted{}~pp unfitted, the strongest single arm anyone constructed against the
deployed policy.

It is nonetheless not an exploit. The excess tracks planted-edge size to within
$\pm$\vsevenInversionTrack{}~pp on all eight rungs of the calibration ladder, which is the signature
of a general strength gain rather than an exploitation instrument. And no linear handicap adds more
than \vsevenHandicapBits{} bits per ask, so within this policy family there is no readability
handicap available for a stochastic policy to purchase. A brief that proposed a stochastic policy
conditional on the inverter succeeding was therefore closed on its antecedent.

\paragraph{A cost figure the record had wrong.} The corpus priced the searching configuration at
\vsevenCorpusCost$\times$ the blueprint. That figure describes the \emph{unrestricted} search, which
is not the frontier point anything was measured at; the endgame-restricted operating point was
re-measured at approximately \vsevenCostRecorded$\times$. \S\ref{sec:results-cost} corrects that
correction in turn.

\subsection{Open adversary generation against the frontier}
\label{sec:development-adversaries}

The constraint on this phase was that the searches must not share a bias, since fifteen runs of one
search constitute one run. Every exploiter search in the corpus before it had been the same search: a
cross-entropy fit of a linear vector maximising win rate from one seed against one target. This phase
constructed six adversary classes across three correlation regimes, and measured channel ceilings
before fitting anything against them.

\textbf{The principal finding is negative.} Nothing constructed exploited the \vsix{} frontier beyond
the detection floor. The best measured exploitability of the deployed policy remains an in-class
figure of \vsevenBestExploit{}~pp $[\vsevenBestExploitLo, \vsevenBestExploitHi]$ over
\vsevenBestExploitN{} games, which is below the class floor of \vsevenPrimaryFloor{}. Seven
independently fitted searches confirmed it and none cleared its class floor.

Every configuration that beat the deployed policy by more than the floor proved on measurement to be
a better policy rather than an exploit: the target's own test-time search at
\vsevenSearchAsExploit{}~pp, and a single out-of-class coordinate at \vsevenOneSwitchArm{}~pp whose
gain survives against opponents sharing none of the target's machinery. The second of those is the
contestation weight of \S\ref{sec:agent-contest}, and it is where this cycle's entire strength gain
originates.

Two further results from this phase bear on \fishseven{}'s design. The declaration-urgency channel
that the phase expected to be its headline has a measured ceiling of \vsevenUrgencyCeiling{}~pp
$[\vsevenUrgencyCeilingLo, \vsevenUrgencyCeilingHi]$ and is dead as an attack, but the same channel is
worth \vsevenUrgencyWorthLo{} to \vsevenUrgencyWorthHi{}~pp to a configuration that simply disables
it: the adversary owns the smaller half of a two-point defect in the target. And the event-count
cliff of \S\ref{sec:agent-urgency} is unreachable by any adversary, while a configuration that removes
the counter inherits a self-play game-length tail of \vsevenCliffTail{} events, which is the
obligation \S\ref{sec:agent-stall} discharges.

\subsection{Candidate mechanisms that were tested and rejected}
\label{sec:development-rejected}

Five candidate mechanisms were developed in isolated branches, gated before being scored, and checked against the side-channel definition by the
harness of \S\ref{sec:evaluation-side} rather
than by inspection. \textbf{One survived}, and it is the configuration change described in
\S\ref{sec:agent} rather than an architecture. The four that were rejected are reported here, each
with the measurement that rejected it, because each closes a standing hypothesis in the project's
design register and because a closed direction does not depreciate.

\subsubsection{The declaration-allocation channel is worth zero}
\label{sec:development-alloc}

The design register's second-ranked hypothesis, at priority \vsevenLonePriority{}, held that
\vsix{}'s entire margin lives in the declaration channel and that roughly two points of it remained.
Its proposed mechanism was that the deployed allocator, a product of per-card marginals, selects a
different allocation from the joint posterior's argmax often enough to matter.

It does not. Over approximately \vsevenLoneDecls{} recorded voluntary declarations, across two banks
and both urgency settings, the joint argmax and the marginal-product argmax select the same allocation
at every one, with the joint score holding a strict opinion on
\vsevenLoneStrictLo\% to \vsevenLoneStrictHi\% of genuinely ambiguous cases. The reason is algebraic
rather than statistical: the sequential joint's first chain-rule factor \emph{is} the marginal, and
the renormalisation that follows reweights the survivors without reordering them.

Two further measurements close the direction rather than merely the mechanism. First, the exact joint
maximiser is a worse decision rule than the deployed approximation on the declaration decision itself:
paired, it corrects \vsevenJointFixed{} declarations and breaks \vsevenJointBroken{}, a net
\vsevenJointWorseNet{}~pp, which converts to approximately \vsevenJointWorseWin{}~pp of win rate. This
is the second independent instance of the methodological rule from \S\ref{sec:background-vsix} paying
off. Second, \vsevenJointFlatLo\% to \vsevenJointFlatHi\% of the error class occupies flat-posterior
states in which the exact object is provably a coin flip, so most of the apparent headroom is an
information limit rather than a policy defect.

The channel is also largely downstream of a defect that \fishseven{} removes for other reasons.
Disabling the urgency escalation eliminates \vsevenLoneUrgencyShare\% of the entire allocation-error
class and drops the oracle ceiling for a perfect allocator from \vsevenLoneCeilingBefore{} to
\vsevenLoneCeilingAfter{}~pp, which is below the detection floor. The consequence is a composition
constraint this paper observes: a declaration-channel gain and the urgency-off gain may never be
added, because they are the same defect counted twice.

\subsubsection{A leaf evaluator fitted 74 times better buys nothing}
\label{sec:development-leaf}

An inherited hypothesis held that the test-time search cannot be judged until its leaf evaluator is
rebuilt, the deployed one being algebraically close to a rescaling of the hit probability. A leaf was
fitted for the purpose, and the statistic that matters, between-candidate $R^2$, which is the only
component the paired deviation rule can consume, improved from \vsevenKoneRsqBefore{} to
\vsevenKoneRsqAfter{}, a \vsevenKoneFold-fold gain replicated out of sample on approximately
\vsevenKoneLeaves{} leaves.

In play it is worth \vsevenKoneInPlay{}~pp against the material control's \vsevenKoneControl{}~pp,
losing on both banks individually. The diagnosis is conversion failure rather than fit failure: the
rule's deviation rate is \vsevenKoneChangeLo{} for the material leaf and \vsevenKoneChangeHi{} for
the retuned one, so at the deployed shrinkage the deviation rule is nearly leaf-invariant. A better
leaf is not the available lever. This reproduces, in a new domain, the abstraction pathology the
poker literature documents \cite{waugh-abstraction}: a model improved on its own metric can leave
play unchanged or worse.

The inherited hypothesis was discharged separately and by re-measurement rather than by the new
evaluator. The single underpowered cell on which it rested, one \vsevenFullGameOldN{}-game cell at
\vsevenFullGameOld{}~pp, re-runs at \vsevenTruncatedEndgame{}~pp replicated in the endgame regime
using \vsix{}'s own leaf.

\subsubsection{Fitting against a per-decision objective produces a worse policy}
\label{sec:development-perdecision}

A further hypothesis held that a per-decision objective is necessary because the scoreboard is too
noisy to resolve small effects. Its precision claim is confirmed and its fitting form is rejected,
and the two halves are independent.

The design effect of deal clustering was measured for the first time in this cycle, and a
one-point-equivalent effect resolves in \vsevenPerDecisionGamesLo{} to \vsevenPerDecisionGamesHi{}
games on the declaration channel against \vsevenScoreboardGames{} on the scoreboard, confirming the
precision claim at \vsevenDesignEffectLo{} to \vsevenDesignEffectHi{} times.

But all three self-oriented per-decision objectives, fitted at matched budget against a per-game
control, produced worse policies: \vsevenKfourSelfDecl{}, \vsevenKfourSelfAsk{} and
\vsevenKfourSelfAlloc{}~pp, each replicated in sign and two of them clearing the floor as confirmed
negatives. Each fit moved its own proxy in the intended direction and lost anyway. The
declaration-oriented fit purchased \vsevenKfourTradeGain{}~pp of declaration accuracy by paying
\vsevenKfourTradeCost{}~pp of ask accuracy, and the ask-oriented fit did the reverse. The two
per-decision channels trade against each other, so a fitter holding \vsevenCoordinates{} coordinates
buys one out of the other. The estimator is worth keeping and the objective is not.

\subsubsection{Nothing learned survives a free random tie-break}
\label{sec:development-tie}

In the \vsix{} mirror \vsevenTieContested\% of contested ask decisions tie bit-for-bit, and the
open question was whether anything can be \emph{selected} within that set rather than merely
randomised. A learned re-ranker
restricted to the tied set was built for the purpose.

Against the deployed policy it scores \vsevenKfiveTieOnly{}~pp, indistinguishable from the free
public-history tie-break's \vsevenKfiveFreeTie{}~pp; head to head against that free tie-break at
\vsevenPhaseTwoCell{} games it is \vsevenKfiveHeadToHead{}~pp $\pm$ \vsevenKfivePm{}, replicated.
Nothing learned survives the free tie-break, and the correct control for any mechanism operating
inside the tied set is the decorrelated policy rather than the deployed one.

The same instrumentation explains why. Setting every candidate's true advantage to exactly zero and
feeding the capture's real standard errors through the engine's real deviation rule reproduces the
search's observed deviation rate to a \vsevenCurseRelErr\% relative error, on a quantity the null was
never fitted to. The search's per-decision selection signal is approximately \vsevenCurseShare\%
winner's curse on its own Monte Carlo draw, and approximately \vsevenCurseTieShare\% within the tied
set, where the rule applies no shrinkage at all.

\subsubsection{Two directions closed on arithmetic}
\label{sec:development-arithmetic}

Closing the entire forced-endgame accuracy gap, re-derived from this cycle's own capture, is worth
\vsevenLthirteenClose{}~pp, which is the register's own figure corrected upward by a factor of two and
still a fiftieth of the detection floor. The single route by which it could have returned, an
adversary driving incidence upward, costs that adversary \vsevenLthirteenAdvCost{}~pp.

\subsection{Selection and freezing}
\label{sec:development-freeze}

\vsevenSelK{} configurations were scored against the deployed policy during development. One was
selected, its specification written to an artifact whose reconstruction is executed rather than
asserted, and the configuration frozen. Nothing about it was changed afterwards, and no fit of any
kind was applied to it. \S\ref{sec:results-selection} quantifies what that selection could be worth
under the null.
